% !PRINTED 622
\weberpage{622}
%<<CONT
führt also die fünf Paare in fünf andere über; und da es nur
zehn Paare von fünf Dingen giebt, so kommen in den beiden
Reihen (3) alle Paare vor.

Man kann nun auf sehr verschiedene Arten halbmetacyklische
Functionen bilden. Die einfachste ist
\weq{4}{u = x_0x_1 + x_1x_2 + x_2x_3 + x_3x_4 + x_4x_0,}
die durch $t$ in
\weq{5}{u' = x_0x_2 + x_2x_4 + x_4x_1 + x_1x_3 + x_3x_0}
übergeht, und $u'$ gehört selbst zu den halbmetacyklischen Functionen.

Ist
\weq{6}{f(x) = x^5 - ax^4 + bx^3 - cx^2 + dx - e = 0}
die Gleichung, deren Wurzeln $x_0, x_1, x_2, x_3, x_4$ sind, so ist also
\weq{7}{u + u' = b.}

Die Function
\weq{8}{y = u - u'}
ist gleichfalls halbmetacyklisch, das Quadrat $y^2$ aber ist vollmetacyklisch
 und also die Wurzel einer Resolvente 6\textsuperscript{ten} Grades,
wie schon \gesperrt{Lagrange} gefunden hat.

Ist $\sqrt{\Delta}$ das Product der zehn Wurzeldifferenzen $(x_0 - x_1)$,
$(x_0 - x_2) \dots$, also $\Delta$ die Discriminante von $f(x)$, so ist
\weq{9}{Y = \dfrac{u - u'}{\sqrt{\Delta}}}

ebenfalls vollmetacyklisch und Wurzel einer Gleichung 6\textsuperscript{ten} Grades $^{1)}$.

Um die conjugirten Functionen zu $u$ zu bilden, bezeichnen
wir mit $M$ die vollmetacyklische, mit $N$ die halbmetacyklische,
mit $A$ die alternirende Gruppe und zerlegen $A$ in die Nebengruppen:
\weq{10}{\begin{aligned}A &= N + N(1,2)(3,4) + Nt(0,1) + Nt(0,2)\\ &+ Nt(0,3) + Nt(0,4).\end{aligned}}

Diese Nebengruppen sind in der That alle von einander
verschieden. Denn wäre z.\,B. $N = N(1,2)(3,4)$, so müsste
$(1,2)(3,4)$ in $N$, also
\[ (1,2)(3,4)t = (1,2)(3,4)(1,2,4,3) = (1,4) \]
in $M$ vorkommen und dies ist nicht möglich, weil in $M$ keine
Permutation auftritt, die zwei Ziffern ungeändert lässt; und in

\par\noindent\rule{4cm}{0.4pt}\par
\noindent $^{1)}$ \gesperrt{Jacobi}, \glqq{}observatiunculae ad theoriam aequationum pertinentes\grqq{},
Crelle's Journ. Bd. 13. \gesperrt{Jacobi's Werke}, Bd. 3. \gesperrt{Cayley}, philos. transactions
 1861, Collected math. papers, Vol. IV, p. 309.
%CONT>>
